%%% Hlavní soubor. Zde se definují základní parametry a odkazuje se na ostatní části. %%%

%% Verze pro jednostranný tisk:
% Okraje: levý 40mm, pravý 25mm, horní a dolní 25mm
% (ale pozor, LaTeX si sám přidává 1in)
\documentclass[12pt,a4paper]{report}
\setlength\textwidth{145mm}
\setlength\textheight{247mm}
\setlength\oddsidemargin{15mm}
\setlength\evensidemargin{15mm}
\setlength\topmargin{0mm}
\setlength\headsep{0mm}
\setlength\headheight{0mm}
% \openright zařídí, aby následující text začínal na pravé straně knihy
\let\openright=\clearpage

\makeatletter
\newcommand\iraggedright{%
  \let\\\@centercr\@rightskip\@flushglue \rightskip\@rightskip
  \leftskip\z@skip}
\makeatother

%% Pokud tiskneme oboustranně:
% \documentclass[12pt,a4paper,twoside,openright]{report}
% \setlength\textwidth{145mm}
% \setlength\textheight{247mm}
% \setlength\oddsidemargin{14.2mm}
% \setlength\evensidemargin{0mm}
% \setlength\topmargin{0mm}
% \setlength\headsep{0mm}
% \setlength\headheight{0mm}
% \let\openright=\cleardoublepage

%% Vytváříme PDF/A-2u
\usepackage[a-2u]{pdfx}

%% Přepneme na českou sazbu a fonty Latin Modern
\usepackage[czech]{babel}
\usepackage{lmodern}
\usepackage[T1]{fontenc}
\usepackage{textcomp}

%% Použité kódování znaků: obvykle latin2, cp1250 nebo utf8:
\usepackage[utf8]{inputenc}

%%% Další užitečné balíčky (jsou součástí běžných distribucí LaTeXu)
\usepackage{amsmath}        % rozšíření pro sazbu matematiky
\usepackage{amsfonts}       % matematické fonty
\usepackage{amsthm}         % sazba vět, definic apod.
\usepackage{bbding}         % balíček s nejrůznějšími symboly
			    % (čtverečky, hvězdičky, tužtičky, nůžtičky, ...)
\usepackage{bm}             % tučné symboly (příkaz \bm)
\usepackage{graphicx}       % vkládání obrázků
\usepackage{fancyvrb}       % vylepšené prostředí pro strojové písmo
\usepackage{indentfirst}    % zavede odsazení 1. odstavce kapitoly

\usepackage[backend=biber]{biblatex}
%\usepackage{natbib}         % zajištuje možnost odkazovat na literaturu
			    % stylem AUTOR (ROK), resp. AUTOR [ČÍSLO]
\usepackage[nottoc]{tocbibind} % zajistí přidání seznamu literatury,
                            % obrázků a tabulek do obsahu
\usepackage{icomma}         % inteligetní čárka v matematickém módu
\usepackage{dcolumn}        % lepší zarovnání sloupců v tabulkách
\usepackage{booktabs}       % lepší vodorovné linky v tabulkách
\usepackage{paralist}       % lepší enumerate a itemize
\usepackage{xcolor}         % barevná sazba

\usepackage{hyperref}
\usepackage[none]{hyphenat}
\iraggedright
\usepackage{acro}	% seznam zkratek
\DeclareAcronym{FPS}{
    short=FPS,
    long=First Person Shooter - střílečka z pohledu první osoby,
}

\DeclareAcronym{RPG}{
    short=RPG,
    long=Role Playing Game - hra na hrdiny,
}
\DeclareAcronym{VR}{
    short=VR,
    long=Virtuální Realita,
}
\DeclareAcronym{LARP}{
    short=LARP,
    long=Live Action Role Playing,
}
\DeclareAcronym{HP}{
    short=HP,
    long=Health Points - body zdraví,
}
\DeclareAcronym{dmg}{
    short=dmg,
    long=Damage - body poškození,
}

\DeclareAcronym{PC}{
    short=PC,
    long=Player Character - postava hráče,
}

\DeclareAcronym{NPC}{
    short=NPC,
    long=Non-player Character - nehráčská (počítačem řízená) postava,
}

\addbibresource{literatura.bib}
%%% Údaje o práci

% Název práce v jazyce práce (přesně podle zadání)
\def\NazevPrace{3D Simulátor šermu s mečem založený na fyzice}

% Název práce v angličtině
\def\NazevPraceEN{3D Physics driven swordfighting simulator}

% Jméno autora
\def\AutorPrace{Jakub Hroník}

% Rok odevzdání
\def\RokOdevzdani{2023}

% Název katedry nebo ústavu, kde byla práce oficiálně zadána
% (dle Organizační struktury MFF UK, případně plný název pracoviště mimo MFF)
\def\Katedra{Katedra softwaru a výuky informatiky}
\def\KatedraEN{Department of Software and Computer Science Education}

% Jedná se o katedru (department) nebo o ústav (institute)?
\def\TypPracoviste{Katedra}
\def\TypPracovisteEN{Department}

% Vedoucí práce: Jméno a příjmení s~tituly
\def\Vedouci{Mgr. Vojtěch Černý}

% Pracoviště vedoucího (opět dle Organizační struktury MFF)
\def\KatedraVedouciho{Katedra softwaru a výuky informatiky}
\def\KatedraVedoucihoEN{Department of Software and Computer Science Education}

% Studijní program a obor
\def\StudijniProgram{Informatika (B1801)}
\def\StudijniObor{IPSS (1801R048)}

% Nepovinné poděkování (vedoucímu práce, konzultantovi, tomu, kdo
% zapůjčil software, literaturu apod.)
\def\Podekovani{%
Poděkování.
}

% Abstrakt (doporučený rozsah cca 80-200 slov; nejedná se o zadání práce)
\def\Abstrakt{%
Boj nablízko s chladnou zbraní najdeme v mnoha videohrách, pouze hrstka z nich se však pokouší o realistickou simulaci, jež by hráči dala svobodu jakkoliv se blížící té, již manipulace s chladnou zbraní umožňuje v reálném světě. Velkou výzvou je návrh schématu ovládání pro klasické počítačové periferie. Ještě méně prozkoumanou oblastí je pak zapojení zbraně do fyzikální simulace herního světa.

Práce implementuje simulátor šermu s jedenapůlručním mečem v enginu Unity, v němž figuruje meč jako plně fyzikálně simulovaný objekt. Rovněž předkládá schéma ovládání pro klávesnici a myš umožňující hráči jemnou kontrolu nad pohybem zbraně. Pro testování hratelnosti implementuje jednoduchého AI protivníka.
Implementace je vytvořena s použitím dobrých programátorských praktik a může komukoliv posloužit jako základ pro akční hru s pokročilým bojovým systémem.
}
\def\AbstraktEN{%
--PLACEHOLDER--
}

% 3 až 5 klíčových slov (doporučeno), každé uzavřeno ve složených závorkách
\def\KlicovaSlova{%
{šerm} {herní fyzika} {počítačová hra} {herní engine Unity}
}
\def\KlicovaSlovaEN{%
{swordfighting} {game physics} {computer game} {Unity engine}
}

%% Balíček hyperref, kterým jdou vyrábět klikací odkazy v PDF,
%% ale hlavně ho používáme k uložení metadat do PDF (včetně obsahu).
%% Většinu nastavítek přednastaví balíček pdfx.
\hypersetup{unicode}
\hypersetup{breaklinks=true}

%% Definice různých užitečných maker (viz popis uvnitř souboru)
%%% Tento soubor obsahuje definice různých užitečných maker a prostředí %%%
%%% Další makra připisujte sem, ať nepřekáží v ostatních souborech.     %%%

%%% Drobné úpravy stylu

% Tato makra přesvědčují mírně ošklivým trikem LaTeX, aby hlavičky kapitol
% sázel příčetněji a nevynechával nad nimi spoustu místa. Směle ignorujte.
\makeatletter
\def\@makechapterhead#1{
  {\parindent \z@ \raggedright \normalfont
   \Huge\bfseries \thechapter. #1
   \par\nobreak
   \vskip 20\p@
}}
\def\@makeschapterhead#1{
  {\parindent \z@ \raggedright \normalfont
   \Huge\bfseries #1
   \par\nobreak
   \vskip 20\p@
}}
\makeatother

% Toto makro definuje kapitolu, která není očíslovaná, ale je uvedena v obsahu.
\def\chapwithtoc#1{
\chapter*{#1}
\addcontentsline{toc}{chapter}{#1}
}

% Trochu volnější nastavení dělení slov, než je default.
\lefthyphenmin=2
\righthyphenmin=2

% Zapne černé "slimáky" na koncích řádků, které přetekly, abychom si
% jich lépe všimli.
\overfullrule=1mm

%%% Makra pro definice, věty, tvrzení, příklady, ... (vyžaduje baliček amsthm)

\theoremstyle{plain}
\newtheorem{veta}{Věta}
\newtheorem{lemma}[veta]{Lemma}
\newtheorem{tvrz}[veta]{Tvrzení}

\theoremstyle{plain}
\newtheorem{definice}{Definice}

\theoremstyle{remark}
\newtheorem*{dusl}{Důsledek}
\newtheorem*{pozn}{Poznámka}
\newtheorem*{prikl}{Příklad}

%%% Prostředí pro důkazy

\newenvironment{dukaz}{
  \par\medskip\noindent
  \textit{Důkaz}.
}{
\newline
\rightline{$\qedsymbol$}
}

%%% Prostředí pro sazbu kódu, případně vstupu/výstupu počítačových
%%% programů. (Vyžaduje balíček fancyvrb -- fancy verbatim.)

\DefineVerbatimEnvironment{code}{Verbatim}{fontsize=\small, frame=single}

%%% Prostor reálných, resp. přirozených čísel
\newcommand{\R}{\mathbb{R}}
\newcommand{\N}{\mathbb{N}}

%%% Užitečné operátory pro statistiku a pravděpodobnost
\DeclareMathOperator{\pr}{\textsf{P}}
\DeclareMathOperator{\E}{\textsf{E}\,}
\DeclareMathOperator{\var}{\textrm{var}}
\DeclareMathOperator{\sd}{\textrm{sd}}

%%% Příkaz pro transpozici vektoru/matice
\newcommand{\T}[1]{#1^\top}

%%% Vychytávky pro matematiku
\newcommand{\goto}{\rightarrow}
\newcommand{\gotop}{\stackrel{P}{\longrightarrow}}
\newcommand{\maon}[1]{o(n^{#1})}
\newcommand{\abs}[1]{\left|{#1}\right|}
\newcommand{\dint}{\int_0^\tau\!\!\int_0^\tau}
\newcommand{\isqr}[1]{\frac{1}{\sqrt{#1}}}

%%% Vychytávky pro tabulky
\newcommand{\pulrad}[1]{\raisebox{1.5ex}[0pt]{#1}}
\newcommand{\mc}[1]{\multicolumn{1}{c}{#1}}


%% Titulní strana a různé povinné informační strany
\begin{document}
%%% Titulní strana práce a další povinné informační strany

%%% Titulní strana práce

\pagestyle{empty}
\hypersetup{pageanchor=false}

\begin{center}

\centerline{\mbox{\includegraphics[width=166mm]{../img/logo-cs.pdf}}}

\vspace{-8mm}
\vfill

{\bf\Large BAKALÁŘSKÁ PRÁCE}

\vfill

{\LARGE\AutorPrace}

\vspace{15mm}

{\LARGE\bfseries\NazevPrace}

\vfill

\Katedra

\vfill

{
\centerline{\vbox{\halign{\hbox to 0.45\hsize{\hfil #}&\hskip 0.5em\parbox[t]{0.45\hsize}{\raggedright #}\cr
Vedoucí bakalářské práce:&\Vedouci \cr
\noalign{\vspace{2mm}}
Studijní program:&\StudijniProgram \cr
\noalign{\vspace{2mm}}
Studijní obor:&\StudijniObor \cr
}}}}

\vfill

% Zde doplňte rok
Praha \RokOdevzdani

\end{center}

\newpage

%%% Následuje vevázaný list -- kopie podepsaného "Zadání bakalářské práce".
%%% Toto zadání NENÍ součástí elektronické verze práce, nescanovat.

%%% Strana s čestným prohlášením k bakalářské práci

\openright
\hypersetup{pageanchor=true}
\pagestyle{plain}
\pagenumbering{roman}
\vglue 0pt plus 1fill

\noindent
Prohlašuji, že jsem tuto bakalářskou práci vypracoval(a) samostatně a výhradně
s~použitím citovaných pramenů, literatury a dalších odborných zdrojů.
Tato práce nebyla využita k získání jiného nebo stejného titulu.

\medskip\noindent
Beru na~vědomí, že se na moji práci vztahují práva a povinnosti vyplývající
ze zákona č. 121/2000 Sb., autorského zákona v~platném znění, zejména skutečnost,
že Univerzita Karlova má právo na~uzavření licenční smlouvy o~užití této
práce jako školního díla podle §60 odst. 1 autorského zákona.

\vspace{10mm}

\hbox{\hbox to 0.5\hsize{%
V \hbox to 6em{\dotfill} dne \hbox to 6em{\dotfill}
\hss}\hbox to 0.5\hsize{\dotfill\quad}}
\smallskip
\hbox{\hbox to 0.5\hsize{}\hbox to 0.5\hsize{\hfil Podpis autora\hfil}}

\vspace{20mm}
\newpage

%%% Poděkování

\openright

\noindent
\Podekovani

\newpage

%%% Povinná informační strana bakalářské práce

\openright

\vbox to 0.5\vsize{
\setlength\parindent{0mm}
\setlength\parskip{5mm}

Název práce:
\NazevPrace

Autor:
\AutorPrace

\TypPracoviste:
\Katedra

Vedoucí bakalářské práce:
\Vedouci, \KatedraVedouciho

Abstrakt:
\Abstrakt

Klíčová slova:
\KlicovaSlova

\vss}\nobreak\vbox to 0.49\vsize{
\setlength\parindent{0mm}
\setlength\parskip{5mm}

Title:
\NazevPraceEN

Author:
\AutorPrace

\TypPracovisteEN:
\KatedraEN

Supervisor:
\Vedouci, \KatedraVedoucihoEN

Abstract:
\AbstraktEN

Keywords:
\KlicovaSlovaEN

\vss}

\newpage

\openright
\pagestyle{plain}
\pagenumbering{arabic}
\setcounter{page}{1}


%%% Strana s automaticky generovaným obsahem bakalářské práce

\tableofcontents

%%% Jednotlivé kapitoly práce jsou pro přehlednost uloženy v samostatných souborech
\chapter*{Úvod}
\addcontentsline{toc}{chapter}{Úvod}

Boj nablízko s chladnou zbraní najdeme v mnoha videohrách. Mezi žánry, které se na něj zaměřují, patří akční adventury, RPG\footnote{\ac{rpg}} a hack-n-slash. Důležitý koncept to je i pro žánr válečných strategických her, doplňující roli pak mívá například v FPS\footnote{\ac{fps}} nebo stealth hrách. Zkrátka, pokud hra ztvárňuje nějaký konflikt řešený násilnou cestou, s vysokou pravděpodobností v ní nějaká chladná zbraň figuruje, ať už jde o světelný meč, středověký řemdih, či židli v hospodské rvačce.

Herní mechaniky, skrz které je boj zblízka realizován, si od počátku éry videoher prošly složitou evolucí a přirozeně každý žánr je adaptuje osobitým způsobem, aby ladily s jeho dalšími specifiky. I tak je ale překvapivě jednoduché popsat určité vlastnosti, které vykrystalizovaly jako společné jádro, jež většinová část her sdílí.

Zpravidla zbraň ve hře funguje jako "černá skříňka" - hráč ji aktivuje stisknutím tlačítka - a postava zaútočí jak sama uzná za vhodné, například přehráním jedné ze seznamu předem připravených animací. Toto prostoduché rozhraní má za výhody přívětivost a snadnou uchopitelnost pro nové hráče a také to, že funguje univerzálně pro víceméně libovolnou zbraň, od jednoručního meče přes halapartnu až po rozbitou lahev od pálenky. Nevýhoda je pak ale, že vzniká výrazný rozdíl mezi tím, jak hráč hraje, a co postava dělá. Hráč tím pádem může mít pocit nedostatečné kontroly a může trpět jeho imerze. Komerčně úspěšných her, které se pokusily hráči dát nad zbraní přímější kontrolu, najdeme v historii pouze několik (\cite{DieByTheSword}, \cite{MountAndBlade}, \cite{KCD}) - podrobněji je rozebereme v první kapitole. Akademický výzkum pak v tomto směru, pokud je autorovi této práce známo, není prakticky žádný, a celkově tedy jde o oblast překvapivě málo prozkoumanou.
  
Chceme-li hráči umožnit přímější kontrolu nad zbraní, mezi očividné výzvy patří návrh schéma ovládání. Klasicky používané periferie jako je počítačová myš či ovladač u herní konzole poskytují dvourozměrnou informaci, chladná zbraň je však obecný objekt v 3D prostoru se třemi stupni volnosti pro pohyb a dalšími třemi pro rotaci, zanedbání jakéhokoliv z nich přinese netriviální důsledky. Je tedy zřejmé, že pomocí klasických periferií není možné dosáhnout dokonale jemného ovládání, které by bylo zároveň pro hráče intuitivní - je třeba nalézt vhodný kompromis. Alternativní periferie, které tento problém nemají, existují (zejména ovladače pro VR\footnote{\ac{vr} - viz kapitolu 1.4.1}), avšak stále nejsou natolik rozšířené, aby s nimi mohl počítat herní mainstream. V dosavadních pokusech herních studií můžeme vidět velmi rozdílné způsoby jak k této výzvě přistupovaly (blíže je představíme v kapitole 1). Každý exceluje v jiných situacích a má jiné očividné nedostatky, z pohledu na ně je zřejmé že tato oblast si vyžaduje další výzkum. 

Jakmile by se nám podařilo pro hráče zajistit jemnou a intuitivní kontrolu nad zbraní, vyplyne nám přirozeně druhá výzva - hráč své kontroly bude chtít využívat. A samozřejmě ne všechny úkony, o které se pokusí, budou mezi těmi, se kterými autor při vytváření hry počítal. Jak zaručíme, že i v takové situaci bude chování hry odpovídat hráčovu očekávání?\footnote{Pomyslete jaké zklamání a pocit zrady musí cítit hráč, který se jen nevinně pokusil podrbat mečem na zádech, načež jeho postava bez varování vystřelila kilometr do oblak a svou smrtí smazala hodinu neuloženého postupu} Zajímavou a dosud neprozkoumanou cestou je pokusit se o plnohodnotnou fyzikální simulaci zbraně. Moderní frameworky pro tvorbu her jsou vybaveny velmi univerzálními subsystémy umožňujícími realistickou fyzikální simulaci pohybu a vzájemné interakce těles, jež na moderním hardwaru nemají problém zpracovávat v reálném čase stovky objektů. Zahrnutím meče do fyzikální simulace můžeme ušetřit díl práce, která by jinak musela být implementována ručně, a zajistíme tím konzistentní chování meče za jakýchkoliv okolností. Navíc tak dosáhneme chování, které již hráči je povědomé z reálného světa, což mu může být nápomocné k rychlému intuitivnímu uchopení systému. V neposlední řadě, meč jako fyzikální objekt může fungovat jako mocné rozhraní pro hráčovu interakci s libovolnými dalšími prvky herního světa, jež jsou rovněž fyzikálně simulované, čímž se otevírá úplně nová oblast emergentních herních mechanik.\footnote{Příklady po krátkém zamyšlení: odrážení střel, srážení hledí protivníkovy helmy, odpalování předmětů na protivníka, páčení víka truhlice, imerzivní tlačítka stisknutelná úderem meče}

Cílem této práce je implementovat kostru jednoduché akční hry z pohledu první osoby - zahrnující hráčem ovladatelnou postavu, prostředí po němž se může pohybovat, velmi jednoduchého počítačem ovládaného protivníka, a zbraň. Zbraní, na kterou se zaměříme, bude \textbf{jedenapůlruční meč}\footnote{\url{https://cs.wikipedia.org/wiki/Dlouh\%C3\%BD_me\%C4\%8D}}. Navrhneme schéma ovládání pro klávesnici a myš, které se neshoduje s žádným dosud používaným, a je koncipované specificky pro naší vybranou zbraň. Budeme při tom usilovat o maximální kontrolu a reakceschopnost hráče při zachování základní uživatelské přívětivosti. Implementaci provedeme v komerčně používaném herním enginu Unity za maximálního využití jeho vestavěného fyzikálního subsystému. Kód bude psán v souladu s dobrými praktikami objektového programování a standarní štábní kulturou enginu Unity (\cite{Unity}), doprovázený podrobnou programátorskou dokumentací a extenzivními testy, aby bylo dalším programátorům umožněno co nejvyšší pohodlí při jeho použítí a rozšiřování. Celý výsledný projekt bude umístěn veřejně pod MIT licencí a zaplní tak prázdný prostor opensource her s jemně ovladatelným bojovým systémem.

\section*{Struktura práce}

V kapitole 1 si blíže vysvětlíme pojmy, se kterými budeme pracovat, uvedeme se do problematiky boje s chladnými zbraněmi ve videohrách, nastíníme příklad typického stavu a několik příkladů, ze kterých chceme čerpat inspiraci. Kapitolu 2 věnujeme stručnému úvodu do použitého herního enginu Unity a podrobnějšímu představení jeho fyzikálního subsystému, na němž tato práce do značné míry stojí. Dále budeme pokračovat kapitolou 3, ve které navrhneme schéma ovládání pro náš jedenapůlruční meč. V kapitole 4 popíšeme celkový návrh a implementaci hry a meče, dopodrobna vylíčíme výzvy, které bylo třeba překonat, a řešení která jsme zvolili. Nakonec v kapitole 5 celkově zhodnotíme návrh našeho systému, rozebereme jeho silné a slabé stránky ve srovnání s existujícími hrami a zamyslíme se nad možnými směry pro navazující práce.
\chapter{Základy}
V této kapitole uvedeme čtenáře do problematiky boje s chladnou zbraní - velmi stručně nastíníme jeho vývoj a metodiku v reálném světě a následně do hloubky rozebereme mechaniky, skrze které ho adaptuje svět akčních videoher, a problémy, které při tom musí řešit.


\section{Chladná zbraň v reálném světě}
V této sekci čtenáře velmi stručně seznámíme s historickým pozadím a základními charakteristikami, které provázejí boj s chladnou zbraní v reálném světě.

\subsection{Průlet historií}

Není nadsázkou když prohlásíme, že chladná zbraň je koncept starý jako lidstvo samo. Archelogické nálezy nasvědčují, že již vzdálený předek moderního člověka Australopithecus používal úderný předmět (pravděpodobně kus dřeva či kost ulovené antilopy) jako \textbf{nástroj k lovu} paviánů \cite{AustralopithecusWeapon}. \textbf{Výhody}, které mu to mohlo přinést, jsou zjevné: zatímco kořist měla možnosti obrany omezené tím, kam dosáhly její vlastní zuby a drápy, lovec si ji mohl držet v bezpečné vzdálenosti a vystavovat vlastní tělo o poznání menšímu nebezpečí odvety.

V průběhu svého vývoje lidský rod zdokonaloval i své zbraně, což přineslo mnoho dalších implikací - ostrý kamenný hrot člověku umožnil zasazovat ránu spolehlivě, a hlubší, než by kdy umožnilo pouhé jeho tělo. Spolehlivost zbraně umožnila lepší \textbf{koordinaci mezi lovci}. Sofistikovaná koordinace člověku otevřela cestu k lovu větších a silnějších druhů zvěře. Nakonec se člověk propracoval na samý vrchol potravního řetězce a zbraň, nástroj lovu, sloužila stále více pro vzájemný \textbf{boj mezi lidmi}. 

Zrod civilizací vznesl na zbraň nové požadavky. Vysoká koncentrace lidí na jednom místě vedla přirozeně ke specializaci profesí, mezi jinými i vojenské. \textbf{Profesionální armáda} čítající tisíce lidí umožňovala (a vyžadovala) do té doby nevídanou míru koordinace - ideálem v takové situaci se ukázalo disponovat širokou škálou zbraní vysoce specializovaných, chladných i střelných, jež mohly být použity ve vzájemné synergii, doplněné válečnými stroji a vhodně vycvičenými zvířaty\footnote{kůň jak známo hrál přední roli}. Profesionální voják měl čas a motivaci svůj typ zbraně pochopit do hloubky, stejně tak vysoce náročná si mohla dovolit být i její výroba a údržba, za níž nesli zodpovědnost rovněž profesionálové ve svém oboru. 
Na druhé straně tu však byl \textbf{běžný obyvatel}, který nepatřil k pravidelnému vojenskému jádru, ale byla-li říše pod útokem, považovalo se za samozřejmost, že pozvedne zbraň na její obranu. Zbraň pro takového člověka vyžadovala především jednoduchost - jak na výrobu, tak na údržbu a použití, zkrátka aby bylo možné v časové tísni krizové situace dostat co nejrychleji co nejvíce odvedenců do bojeschopného stavu a udržet je v něm.

Mocenské soupeření států ústí ve zběsilý a nikdy nekončící \textbf{závod ve vývoji účinnějších zbraní a metod jejich použití}. V jednom okamžiku dominuje Blízkému východu vozataj, v následujícím jej poráží Makedonská falanga, nad ní předvede svou superioritu Římský systém manipulů, ten vzápětí Římané sami prohlašují za zastaralý, avšak o půl tisíciletí později stejně jejich říši rozvrací hordy Hunských jezdců... chladná zbraň, ať už v rukou jezdce či pěšího vojáka, zůstává po většinu dějin dominantním prvkem na bojišti, se střelnými zbraněmi hrajícími významnou podpůrnou roli. Až s rozmachem zbraní palných se tato dynamika obrací a velmi pozvolně se dobíráme k modernímu vojenství. V současné době se chladná zbraň považuje převážně za překonanou - uplatnění pro ní stále existuje např. v kontextu pořádkových složek, avšak i pro ty plní úlohu spíše doplňkovou. Mimo oficiální kruhy pořád hraje nezanedbatelnou roli, avšak to je především pro její \textbf{triviální dostupnost} v porovnání s palnou zbraní.

Tradice chladné zbraně však stále žije v civilních komunitách dedikovaných zachování historie a kulturního odkazu. Mezi významné patří japonské umění \textbf{Kendó} ("cesta meče") vycházející ze samurajské tradice, či evropská komunita \textbf{historického šermu}, jež vychází z dochovaného učení středověkých mistrů. Rovněž je zde moderní \textbf{sportovní šerm}, jež přímo navazuje na šermířskou tradici ranného novověku. V posledních letech tyto vlivy více pronikají i do běžných volnočasových aktivit - ve Střední Evropě například je stále oblíbenější fenomén \textbf{LARP}\footnote{\ac{LARP}}, jehož jedna z podob - tzv. "dřevárna" - znamená hromadnou akci, na které se desítky až stovky lidí v tématických kostýmech a vyzbrojení zpravidla dřevěnými, molitanem měkčenými replikami zbraní, střetnou v bitvě. Pravidla boje přirozeně musí zaručit bezpečí účastníků, avšak stále při zachování autentického zážitku z boje.  

Vzhledem k charakteru a rozšířenosti výše zmíněných fenoménů lze očekávat, že jejich zastánci mají v nemalé míře zastoupení i ve videoherní komunitě. Přirozeně takoví lidé mohou mít zájem o hru, jež jim umožní jejich oblíbenou činnost napodobit, ale bez mnohých nepohodlí, jež ji provázejí v reálném světě. Pro hráče, který s bojem v reálném světě žádnou zkušenost nemá, může zase videohra být cestou, jak tuto mezeru ve svých zkušenostech částečně zaplnit a získat například hlubší porozumění k historii.

\subsection{Průběh boje}
Nyní si představíme typický průběh boje vedeného s chladnými zbraněmi, a vypíchneme na něm znaky, které by videohra měla dokázat vystihnout, pokud se pokouší o realistický bojový systém.
\bigbreak

Cílem souboje je v nejjednodušším případě protivníka \textbf{zabít, vážně zranit či ho jinak vyřadit z bojeschopného stavu}, čehož může být dosaženo libovolnými prostředky, a ideálně zabránit, aby při tom stejný úděl potkal i vás. Není vyloučeno vítězství vyčerpáním protivníka či využití okolních předmětů a znalosti prostředí (např. lstivým vlákáním protivníka do bažiny).

Tato absolutní svoboda však může být omezena morálními zásadami účastníků\cite{HistoryOfSurrender} či předem dohodnutými pravidly souboje. Velmi zde záleží na \textbf{okolnostech}, při kterých k boji dochází:
\begin{itemize}
    \item Např. při \textbf{obraně domácnosti} před vetřelcem platí popis výše
    \item \textbf{Voják v bitvě} koná v mezích, jež mu stanovuje vojenská disciplína a rozkazy velitele
    \item \textbf{Účastník turnaje} přirozeně nesmí opustit vyhrazené bojiště či využívat pomoci z publika
    \item Ve \textbf{cvičném souboji} si účastníci logicky dávají pozor, aby jeden druhého zbytečně nezabil či nezmrzačil
    \item V duelu na obranu cti se účastníci mohli \textbf{předem dohodnout}, že nechtějí jeden druhého zabít, a např. bojovat "na první krev"
\end{itemize}
Mezi znaky zkušeného šermíře patří schopnost být si vědom těchto mantinelů a v jejich rámci jednat kreativním a účinným způsobem.

Extrémní případ těchto omezení vidíme v moderních šermířských komunitách. U těch najdeme velmi \textbf{striktní pravidla} a vysoké nároky na bojovou výstroj s cílem zamezit jakýmkoliv vážným zraněním a ztrátám na životech, jelikož ty přirozeně nejsou cílem jejich snažení. 

Avšak pokud nedochází k žádným reálným zraněním, \textbf{jak lze rozhodnout, kdo vyhrál?} Zjevnou cestou je například systém zásahových zón a počítání životů, které se bojovníkovi v závislosti na zasažené zóně odečítají. Takový systém však nezaručuje realistický dojem z boje. Mezi plody snah o realističtější zážitek patří např. v české LARPové scéně oblíbený šatrh\cite{Satrh}, založený na důvěře v účastníky, že podle způsobu, jakým byli zasaženi, odhadnou a zahrají realistickou reakci, či buhurt, kde zkrátka vítězí ten, kdo nepadne vyčerpáním. Je tedy vidět, že i komunity, jejichž nejvyšší prioritou při boji je záruka bezpečí, jsou ochotné vyvinout značnou snahu, aby byl zážitek z boje realistický.

\subsection{Boj jeden na jednoho}
Jak tedy probíhá duel mezi dvěma šermíři? Hloubkový rozbor zohledňující historické i moderní techniky, rozdílné úrovně zkušenosti, zdatnosti a motivace obou zápasníků, veškeré obvyklé zbraně, kterými by mohli být vybaveni, prostředí a vnější okolnosti, které v souboji mohou hrát podstatnou roli, není možné provést v rozsahu bakalářské práce, natož jedné její kapitoly. Pro ten tedy čtenáře odkáži na specializovanou literaturu \cite{KunstDesFechtens} \cite{FightingWithTheGermanLongsword} \cite{ModernHEMA}.  

Zde si pouze nastíníme pár základních faktorů, které při souboji hrají důležitou roli.

\textbf{Obratnost zbraně} - schopnost rychle měnit její trajektorii, reagovat na změny situace (např. vyblokovat nečekaný úder) či změny situace sám vyvolat (např. na poslední chvíli zastavit fingovaný úder a stočit čepel na jiné místo, které protivník nechal nechráněné). Mezi faktory, kterými je determinována, patří velikost a hmotnost zbraně, její tvar, těžiště a způsob úchopu\footnote{Čím jsou ruce dál od sebe, tím silnější je princip páky.}; v neposlední řadě také síla a tělesné rozpoložení šermíře.

\textbf{Dosah, pohyb, terén} - V okamžiku, kdy vám protivník není schopen ublížit a vy jemu ano, můžete si dovolit být libovolně agresivní\footnote{Díky tomu dlouhé zbraně jako kopí mohou být velmi účinné i v rukou nezkušeného odvedence.}. Je-li váš protivník v takovéto výhodě, jedinou možností je se k němu přiblížit - pak se dostanete do výhody vy, jelikož zbraně s dlouhým dosahem zpravidla trpí na nízkou obratnost při boji zblízka\footnote{Nedejte protivníkovi čas tasit poboční zbraň.}. Vaše možnosti, jak se k němu přiblížit, přirozeně velmi závisí na terénu\footnote{Když stojíte pod hradbami a protivník bodá kopím z ochozu nad vámi, jste v pěkné kaši.}. 

Máte-li oba stejný dosah, pohyb se však nestává o nic méně důležitým. Správná technika pohybu pomáhá šermíři držet pod kontrolou jeho těžiště a nenechat se rozhodit, když klopýtne, či na něj protivník agresivně naběhne; případně mu umožní bleskově se přiblížit a využít nechráněnou oblast, kterou mu protivník nabídl svým nemotorným pohybem. 

\textbf{Zbroj}, kterou má protivník na sobě, může citelně omezit vaše možnosti co do způsobů, jak mu uštědřit zranění. Je-li například v plné plátové zbroji pozdního středověku, pokoušet se do ní sekat mečem je prakticky zbytečné. Vaše možnosti se citelně zúží - buď se můžete pokusit bodnout do některé ze skulin ve zbroji, nebo meč vyměnit za palcát či válečné kladivo\footnote{Případně lze meč chytit za čepel a záštitou zasuplovat hlavici kladiva - viz \cite{FightingWithTheGermanLongsword}} a protivníkovi tupými údery zlámat kosti. Protivník si vaší situace bude vědom, a o to jednodušší pro něj bude předvídat vaše akce.

Znakem každého dobrého šermíře je schopnost vcítit se do pozice svého protivníka a \textbf{odhadovat jeho budoucí akce} dříve, než se stanou. Nezkušený šermíř bude každý svůj úder provázet výrazným nápřahem, pro jeho protivníka pak není obtížné úder vyblokovat a přitom se rovnou k protivníkovi přiblížit a seknout po místě, které ví, že při úderu protivník nechal nechráněné. Zkušený šermíř pak ani nepotřebuje vidět výrazný nápřah, naučí se protivníkovu akci podvědomě vycítit z drobných náznaků (pohybů svalů, očí,...). Samozřejmě je možné předvádět údery fingovaně a čekat, že se při protiútoku otevře protivník vám. Pak záleží na jeho úsudku, zda úmysl prokoukne a provede fingovaný protiútok - stupňům falše se meze nekladou. 

\subsection{Hromadný boj a boj proti přesile}
Souboj jeden na jednoho však může vypadat jako téměř laboratorní podmínky vedle chaotických situací, v rámci kterých v historii běžně k boji docházelo.   

Často se střetávaly větší, více či méně \textbf{organizované skupiny}, případně celé armády čítající mnoho takových pevně ustanovených jednotek. V takové situaci je na místě, aby bojovníkova individualita ustoupila do pozadí - rozhodujícím faktorem je až na výjimky taktické jednání a sladěnost celku. 

Jedním ze základních požadavků, které jsou na vojáka kladeny, je schopnost \textbf{držení formace}. Stojíte-li jako semknutá řada, štít na štít, meče připravené opětovat jakýkoliv úder, který by protivník zkusil na kamaráda vedle, je mnohem těžší prorazit vašimi řadami a dostat se k lučištníkům, kteří nepříteli způsobují těžké ztráty. 

Při skupinové šarvátce hraje zpravidla rozhodující roli taktická \textbf{koordinace} spolubojovníků. Nepřítel, který vám překvapivě vpadl do zad, je mnohem nepříjemnější záležitost, než nepřítel, který postává opodál a čeká, až na něj přijde řada.  

Pro dosažení koordinace je však nutná \textbf{komunikace}. Dohodnout se při šarvátce na postupu, v omezeném čase a aniž by plán byl zachycen protivníkem, vskutku není triviální problém. V kontextu větší bitvy pak důležitost i obtížnost komunikace nabývá zcela nového rozměru a jde o těžký problém, který až do vynálezu moderní vysílačky nebyl uspokojivě vyřešen. Vojenská taktika si vyžaduje flexibilní reakce na změny situace na bojišti, informace musí proudit mezi vojenskými jednotkami a velitelstvím, které jsou často netriviálně fyzicky vzdálené. Jednou možností, jak informaci předat, je vyslání posla - ten je schopen nést detailní informaci, je však zranitelný a pomalý - než dorazí, může již jeho informace být neaktuální. Protipólem pak je broadcast předem dohodnutých signálů, napříkald prostřednictvím trubače - tato cesta je rychlá a v určitém doslechovém okruhu vcelku spolehlivá, avšak signály musí být předem dohodnuté a variant nemůže být mnoho, metoda tedy postrádá flexibilitu a detailnost informace. Jak vidno, jde o zásadní součást vedení bitvy, přesto však vidíme mnoho videoher, které ji zcela zanedbávají.
\bigbreak
Dojde-li k boji \textbf{jediného bojovníka proti přesile}, vše výše uvedené stále platí, akorát očividně kolosální výhody plynoucí z koordinovaného postupu jsou dopřány pouze jedné ze stran. I pro velmi zkušeného šermíře bojujícího proti několika nezkušeným, často jediná šance, jak z takové situace vyváznout, je protivníky rozdělit (kreativním použitím pohybu a znalosti terénu) a zdolat jednoho po druhém rychle, než mu ostatní přijdou na pomoc.

\subsection*{Shrnutí}
Seznámili jsme čtenáře s historickým pozadím a základními prvky, které charakterizují boj s chladnou zbraní v reálném světě. Nyní by měl být připraven pokročit dál a provést srovnání s přístupem, jakým tuto problematiku adaptuje svět videoher.

\clearpage
%------------------------------------------------------------------------------------------------------------------------------------------------------------------------------------------------------------------------------------------------------------%
 % xxxxxxxxxxxxxxxxxxxxxxxxxxxxxxxxxxxxxxxxxxxxxxxxxxxxxxxxxxxxxxxxxxxxxxxxxxxxxxxxxxxxxxxxxxxxxxxxxxxxxxxxxxxxxxxxxxxxxxxxxxxxxxxxxxxxxxxxxxxxxxxxxxxxxxxxxxxxxxxxxxxxxxxxxxxxxxxxxxxxxxxxxxxxxxxxxxxxxxxxxxxxxxxxxxxxxxxxxxxxxxxxxxxxxxxxxxxxxxxxxxxxxxxx %
%------------------------------------------------------------------------------------------------------------------------------------------------------------------------------------------------------------------------------------------------------------%


\section{Chladná zbraň v akčních videohrách}

Nyní, když máme rámcovou představu, jak boj s chladnou zbraní vypadá v reálném světě, nastíníme si, jakým způsobem bývá typicky adaptován v akčních videohrách - jaké jeho aspekty se podařilo v tomto universu věrně zpodobnit, jaké byly ignorovány, a jaké interpretovány s řádnou dávkou umělecké kreativity.

\subsection{Obecný přehled}

Účinkujícími ve hře jsou \textbf{herní postavy}. Ty mohou být ovládané buď hráčem (takovou postavu nazveme \textbf{\acs{PC}}\footnote{\Acl{PC}}), nebo počítačovým algoritmem (tu označíme jako \textbf{\acs{NPC}}\footnote{\Acl{NPC}}), a lze je popsat nějakým jejich \textbf{vnitřním stavem}. Nyní si představíme základní veličiny, které jsou typickou součástí vnitřního stavu postavy ve všech akčních hrách bez ohledu na to, zda jde o hru zaměřenou na boj se střelnou zbraní, chladnou zbraní, či boj jakýkoliv jiný:

\begin{itemize}
    \item \textbf{\ac{HP}}... Číslo udávající zdravotní stav a životaschopnost postavy. \textbf{Pokud klesne na nulu, postava umírá.} K jeho snížení může dojít například úderem protivníkovy zbraně či pádem z výšky a takto odebraným bodům HP říkáme \textbf{\acs{dmg}}\footnote{\Acl{dmg}}. Hráči je stav jeho HP obvykle v reálném čase explicitně komunikován skrze ukazatel na obrazovce. V průběhu boje typicky nejsou doplňovány\footnote{Popřípadě prostředky, skrze které je lze doplnit, jsou vzácné.} a postava musí vystačit s tím, co má - jde tedy o zdroj vybízející k dlouhodobějšímu plánování.
    \item \textbf{Výdrž}... Volitelný avšak častý prvek, číslo reprezentující akutní schopnost postavy vynakládat fyzické úsilí. Spotřebováváno chvatným pohybem a prováděním akcí, v průběhu boje zpravidla dochází k jeho rychlému opakovanému spotřebování a regeneraci.
    \item \textbf{Inventář}... List dalších předmětů, které postava nese s sebou. Může jít o peníze, léčivé lektvary, munici, náhradní zbraně, magické svitky, mapy a cokoliv dalšího v závislosti na hře. 
\end{itemize}

V závislosti na vstupu, který získává od hráče či od řidícího algoritmu, se postava pohybuje po herním světě, interaguje s objekty herního světa a s obsahem svého inventáře a používá svou zbraň. Na použití zbraně se nyní zaměříme.

\textbf{Ovládání zbraně} bývá typicky omezené na stisknutí tlačítka - \textbf{Zaútoč!} - načež postava sama provede útok tak, jak uzná za vhodné - obvykle přehráním jedné ze seznamu animací, které pro danou kombinaci postavy a zbraně předem připravil autor hry.

Implikace, které takto jednoduché rozhraní přináší \textbf{pro hráče}, jsou následující:
\begin{itemize}
    \item \textbf{Intuitivnost} - Hráč, který poprvé spustí hru a zkusí náhodně mačkat tlačítka, rychle ovládání pochopí a začne být ve hře použitelný - žádná nutnost procházet úmorným tutoriálem.
    \item \textbf{Uniformní rozhraní pro všechny zbraně} - Tímto způsobem lze ovládat libovolnou zbraň, ať už jde o dvoumetrovou halapartnu či pěsti hráčovy postavy. Hráč tedy není přehlcen množstvím jednoúčelových herních systémů a z výhod, které získá svým zdokonalováním v použití bojového systému, může těžit po celou hru.
    \item \textbf{Nízká specifičnost} - Rozhraní hráči neumožňuje těžit ze silných stránek žádné konkrétní zbraně. Hra mu skrze bojový systém není schopna poskytnout vzdělávací lekci ohledně vlastností zbraně v reálném světě. Pokud mu tyto vlastnosti již jsou známé, může utrpět jeho imerze.
    \item \textbf{Limitovaná předvídatelnost} - Vzhledem k tomu, že akci, která bude vykonána, nevybírá hráč, ale jeho herní postava, vnáší se tím do hry prvek náhody, který má potenciál hru emocionálně okořenit.\footnote{Pokud však výběr akce probíhá nějakým deterministickým způsobem, hráč v něm brzy intuitivně vypozoruje určité vzory a naučí se akce své postavy do jisté míry předvídat a kalkulovat s nimi při plánování svého postupu - tím může bojový systém získat na hloubce.}
\end{itemize}

A co důsledky, které tento přístup přináší z pohledu \textbf{tvůrce hry}?:
\begin{itemize}
    \item \textbf{Předvídatelnost} - Tvůrce hry má přesnou kontrolu nad tím, jaké akce může jakákoliv postava potenciálně vykonat. Tím se zmenšuje prostor pro výskyt neošetřených okrajových případů.
    \item \textbf{Uniformní rozhraní pro všechny zbraně} - Herní logika, která umí pracovat s jednou zbraní, umí pracovat s každou zbraní. To může dopomoci k celkově čistému a kvalitnímu kódu a značně se tím snižuje námaha s přidáváním dalších zbraní (či jejich procedurální generací za běhu).
    \item \textbf{Nízká specifičnost} - Není možné uspokojivě ztvárnit veškerá zajímavá specifika konkrétní zbraně. Všechny zbraně se musejí chovat do jisté míry podobně. Často je tak zbraň smrštěna do pouhé tabulky statů\footnote{Povrchní veličiny typu dosah, rychlost úderu, míra způsobeného zranění}.
\end{itemize}

Vidíme, že takovéto pojetí boje se zbraní má mnohá pozitiva, stále je však očividné, že \textbf{jde o systém jednoduchý, který sám o sobě nemůže sloužit jako nosná herní mechanika v akční videohře} - oblasti, kde cílová hráčská skupina očekává určitou míru komplexity a hloubky.

Hloubka bývá tomuto systému zpravidla dodávána výběrem mezi několika různými typy útoku (každý efektivní v jiné situaci), akcí pro blokování nepřátelských útoků, důrazem na časování (krátkodobé omráčení nepřítele, navazování úderů, odměna za vyblokování nepřítele v přesně správném okamžiku,...) a \textbf{silným propojením s dalšími mechanikami}, jakými je v první řadě pohyb, dále různé lektvary, magie apod. v závislosti na hře.

Toto si nyní názorně předvedeme na příkladě komerčně úspěšné hry z nedávných let.

\subsection{Zaklínač 3 - vzorový příklad}
Zaklínač 3: Divoký Hon \cite{Wiedzmin3} je akční \acs{RPG} hra, která byla vydána r. 2015 polským studiem CD Projekt. Mezi její hlavní důrazy patří \textbf{otevřený fantasy svět}, propracovaný \textbf{nelineární příběh}\footnote{Hra vychází z knižní série Zaklínač Andrzeje Sapkowského \cite{SapkowskiZaklinac}} a v neposlední řadě také dynamický \textbf{bojový systém}.

Hráč se zde ujímá role Geralta z Rivie - zaklínače, \textbf{profesionálního lovce monster}. Ten je kromě svých dvou mečů\footnote{stříbrného pro boj s nestvůrami a železného proti lidem} vybaven flexibilními možnostmi pohybu, jednoduchými kouzly (tzv. Zaklínačská Znamení) a škálou ručně vařených lektvarů. Na své cestě bojuje s lidmi, duchy, upíry, bazilišky i nepopsatelnými obludami z nočních můr.
Jak tedy vypadají mechaniky, skrze které bylo toto vše umožněno?:
\bigbreak

Základem jsou \textbf{dva typy úderů - rychlý a silný}. Jak vyplývá z názvu, rychlý úder je rychlejší, oproti tomu animace silného úderu trvá mírně delší čas, avšak zranění způsobené protivníkovi je také adekvátně zvýšené. K čemu je to dobré? Pro mnoho protivníků platí, že v okamžiku, kdy jim je učiněno nezanedbatelné zranění, přeruší akci, kterou se chystali provádět, a na zlomek času zakolísají. S rychlým úderem je hráč schopen toto zakolísání vyvolat častěji - s dobrým časováním se lze například dostat do cyklu, kdy lehkooděného protivníka hráč udolá jedním úderem za druhým aniž ten se zmohl udělat cokoliv proti vám. Toto stejné zakolísání však může potkat i hráče. Hráč se tedy neustále nachází v časovém tlaku, který ho vybízí kalkulovat - např. zda stihne dva rychlé údery, či raději jen jeden silný, než ho nějaký nepřítel přeruší. Dalším důvodem proč používat silný útok jsou těžkoodění nepřátelé, kterým rychlý útok udílí zranění velmi redukované. 

Drobným doplňkem je akce \textbf{blokování}. Bojujete-li proti lidem, držení tlačítka blokování vás činí prakticky imunním proti běžným útokům mířícím na vás zepředu, za tu velmi drobnou cenu, že sám nemůžete útočit a pohybovat se po bojišti lze jen velmi pomalu. Zajímavější je \textbf{perfektní blokování} - pokud zahájíte blokování přesně v okamžiku, kdy na vás protivník vede úder, provedete odvetný protiúder, který protivníka na pár sekund omráčí. Tím je do boje vnesen další prvek časování. Rovněž je třeba dávat pozor na nepřátele, kteří blokují proti vám.

Proti chaotickým útokům vedeným změtí zubů a drápů nestvůry však blokování mečem ztrácí smysl. Tehdy přichází na řadu možnosti úskoků a kotoulů, které nabízí \textbf{systém pohybu}. Kromě nich samozřejmě hra nabízí i klasickou chůzi a sprint. Úskok je velmi drobný posun do strany, který nespotřebovává výdrž, pouze má interní cooldown (několik desetin sekundy) zabraňující spamování. Cílem je umožnit hráči tak akorát uskočit z oblasti dopadu nepřítelovy zbraně. Pro zajištění, že tento záměr bude vždy naplněn, je postavě na malý zlomek sekundy poskytnuta celková nezranitelnost. Kotoul je pak rámcově obdobný s tím rozdílem, že spotřebovává výdrž a na oplátku ho lze použít k rychlému překonání velké vzdálenosti. Hra klade velký důraz na hromadné boje, kde hráč musí být schopen prioritizovat cíle a rychle využívat objevivších se zranitelností jednotlivých nepřátel. Vysoká mobilita, kterou kotoul poskytuje, je tedy více než nápomocná.  

Zaklínačská Znamení do boje přinášejí další vrstvu komplexity. Jde o \textbf{5 kouzel}, která spotřebovávají velkou porci výdrže a na oplátku zaklínači mocně pomohou - např. může jít o ochranné silové pole, poryv větru, který má šanci omráčit zasažené protivníky, či zmatení, které přiměje nepřítele bojovat na vaší straně.

Výčet Znamení je pevně stanoven, avšak klasicky pro žánr RPG, hra zahrnuje \textbf{prvky progrese}, pomocí kterých lze Znamení v průběhu hry vylepšovat a modifikovat. Stejné platí i pro boj s mečem - rychlý i silný úder mají každý své vlastní stromy dovedností, které jsou schopny zajímavými způsoby ovlivnit gameplay. 

Poslední částí zaklínačova repertoáru, kterou zmíníme, jsou lektvary a oleje. Jde o \textbf{konzumovatelné předměty}, které hráč v průběhu hry sám vyrábí, a po jejich konzumaci (požití lektvaru postavou / namazání oleje na meč) mu je na určitý časový limit poskytnut nějaký bonus (př. regenerace zdraví, bonusové zranění proti upírům apod.). Jde tedy o zdroje, se kterými hráč hospodaří v dlouhodobějším horizontu, ale použity v kritickém okamžiku mohou zvrátit výsledek boje. 

Esenciální součástí každé hry s bojovými prvky, jsou také \textbf{nepřátelé}. V této oblasti byl Zaklínač 3 k hráči velmi štědrý. Najdeme zde obyčejné lidské nepřátele (bandity, nepřátelské vojáky apod.), jejichž způsob boje sdílí mnoho prvků se zaklínačovým - např. jsou schopni blokovat jeho údery. Dále pak širokou paletu monster všech velikostí a tvarů, pozemní i schopné létat vzduchem, vybuchující, duchy zranitelné pouze specifickými cestami, tvory vládnoucí magií - každý nepřítel má své osobité silné a slabé stránky a vybízí hráče k jinému stylu hraní.

\bigbreak

Vidíme, že bojový systém Zaklínače 3 je zjevně postaven kolem kostry, kterou jsme si načrtli v předchozí kapitole, zároveň jde o systém mající značnou hloubku. Tato hloubka je získána z velké části důrazem na celkový hráčův přehled o situaci na bojišti, škálou osobitých nepřátel, důrazem na časování, a pak také propojením s dalšími systémy - pohybem, magií a vařením lektvarů. 

Hráči je poskytnut prostor k budování dovednosti a k důmyslnému taktizování jak v krátkodobém, tak v dlouhodobém horizontu. Celkové kvality bojového systému ostatně potvrzuje značný komerční úspěch, kterého se hře po vydání dostalo.

Mezi možné příčiny úspěchu hry můžeme zařadit také pozorování, že \textbf{hloubka systému je volitelná}. Rekreační hráč, který hraje převážně kvůli příběhu, si ji může dovolit víceméně ignorovat. Stačí hru nastavit na nízkou obtížnost a většinu nepřátel není těžké udolat pouze s povrchní znalostí základních mechanik. Není odvážné tvrdit, že to může mít pozitivní vliv na velikost cílového hráčského publika. 

\subsection{Chladná zbraň jako doplňkový prvek}

Chladnou zbraň najdeme také v mnoha hrách, ve kterých není stěžejním prvkem. Například ve většině her žánru \textbf{\acs{FPS}} najdeme kromě širokého výběru palných zbraní i nějaký nůž jako poslední zálohu když dojde munice, v žánru \textbf{stealth} často vidíme chladnou zbraň jako prostředek k tiché likvidaci nepřátel a takto bychom mohli pokračovat. Mnohdy taková zbraň odpovídá popisu z 1.2.1 - cílem není, aby chladná zbraň sama o sobě byla zajímavou mechanikou, naopak \textbf{tvoří pouze jeden dílek} ve skládačce herního systému a pro hráče musí především být jednoduše uchopitelná. 

Přesto najdeme mnoho takových zbraní, které se staly \textbf{ikonickým symbolem} své hry - zmínit můžeme například hru DOOM \cite{DOOM1993} a motorovou pilu, či páčidlo ze série Half-Life \cite{HalfLife}.


\section{Hry s větší kontrolou}

Nyní víme, jakým způsobem je typicky chladná zbraň v akčních hrách ztvárňována, známe pozitiva, která tato cesta přináší, a rovněž některé triky, jak jejím prostřednictvím dosáhnout poutavého bojového systému. Co nás ale čeká, kdybychom chtěli sejít z vyšlapané cesty a nabídnout hráči nad jeho zbraní větší míru kontroly?

\subsection{Benefity a výzvy}

Neprve se zamysleme, \textbf{proč by hráč o větší kontrolu mohl mít zájem}:
\begin{itemize}
    \item \textbf{Je to poučné} - Jedním z hlavních cílů videoher celkově je zprostředkování zážitku, který je pro hráče z různých důvodů nedosažitelný v reálném světě. Skrze hru, která mu zprostředkuje přímou kontrolu nad mečem, může hráč získat větší pochopení k problematice historického vedení boje.
    \item \textbf{Šerm ho baví v reálném světě} - Naopak, hráč může být aktivním šermířem, avšak chce si zážitek z boje užít ve zjednodušené podobě i v pohodlí domova.
    \item \textbf{Nemá rád když ho hra omezuje} - Nemalá skupina hráčů si ráda hledá vlastní cestu hrou a vyžaduje od hry plnou svobodu v tom, co jim je umožněno dělat. 
    \item \textbf{Chce zkusit něco nového} - Mnoho hráčů se již zkrátka nudí z donekonečna omýlaného klišé, jakým se v jejich očích boj s chladnou zbraní stal - líbí se jim téma, ale zároveň ho chtějí zkusit pojaté neotřelým způsobem.
\end{itemize}

Dále, jaké benefity by takový přístup mohl přinést \textbf{tvůrci hry}:
\begin{itemize}
    \item \textbf{Emergentní gameplay} - Tvůrci hry stačí, aby zprovoznil robustní systém, který hráči umožní obecný pohyb meče. Najít zajímavou metodiku jak systém použít, je již úkolem hráčské komunity. Takto může dojít k živelnému objevení mechanik, které tvůrce hry nikdy ani neplánoval, čímž hra může být značně obohacena.  
    \item \textbf{Nástroj interakce s okolním prostředím} - Pohybem meče v 3D prostoru se dá zprostředkovat široká škála interakcí s okolním prostředím, které by jinak musely být implementovány odděleně (zmíníme např. mačkání tlačítek či páčení víka truhlice). Také se tím otevírá cesta pro zcela nové interakce, které se s použitím meče jeví přirozené, bez něj by však bylo těžké najít cestu, jak je do hry zakomponovat (např. odrážení střel, srážení hledí protivníkovy přilby, odpalování předmětů na protivníka apod) .   
\end{itemize}

Vidíme, že motivace k umožnění větší kontroly je silně nezanedbatelná. Očividně tedy takovým pokusům musí stát v cestě výrazné překážky vzhledem k tomu, že víme, že se jim dosud nepodařilo ovládnout herní mainstream.

\textbf{Jaké očividné problémy je tedy třeba vyřešit?}
\bigbreak
Prvním je \textbf{ovládání}. Meč je obecný objekt v 3D prostoru, má tři stupně volnosti pro polohu a další tři pro rotaci. Nemluvě o dalších vlastnostech - např. jak pevně ho šermíř drží. Pro žádnou z těchto veličin se nenabízí očividný způsob, jak ji zanedbat či odvozodit z těch ostatních. Abychom mohli mečem sekat, bodat a blokovat údery, potřebujeme pohybovat rukojetí a natáčet čepel, obojí po všech třech osách. Rotace meče kolem jeho středové osy by se zdála jako vhodný kandidát k implicitnímu odvozování, avšak např. k odpalování projektilů je velmi vhodné moci zajistit, že do nich udeříme plohou stranou meče, místo abychom je rozsekli ostrou hranou čepele.

Oproti tomu \textbf{počítačové periferie}, které se běžně používají jako zdroj hráčského vstupu - tedy myš a klávesnice, popřípadě ovladače u herní konzole - poskytují informaci nanejvýše dvourozměrnou (horizontální a vertikální osa). Neexistuje tedy přímočarý a intuitivní způsob, jak na ně pohyb meče namapovat - je nutné učinit kompromis a nějakým způsobem volnost pohybu meče omezit. Alternativní metody ovládání, které tento problém nemají, existují, a stručně se na ně podíváme později v této kapitole, jejich rozšířenost mezi hráči však není taková, aby bylo rozumné spoléhat, že jimi naše cílová skupina disponuje.
\bigbreak
Pokud se nám podaří problém ovládání uspokojivě vyřešit, narazíme na další v oblasti \textbf{hratelnosti}. Důležitou součástí dynamiky boje s chladnou zbraní je předvídání protivníkových úderů a reagování na ně ještě dříve, než se stanou. V reálném světě je šermíř schopen pozorovat drobné pohyby svalů a očí protivníka, které v tomto ohledu mnohé napoví. V počítačové hře samozřejmě je hráčova schopnost vnímat takovéto detaily velmi omezená. Tento problém je ve hrách typicky řešen tak, že animace úderů zahrnují jako počáteční fázi výrazný nápřah. Jak však zajistíme takovéto výrazné telegrafování útoků, pokud má nepřítel nad mečem jemnou kontrolu a může z principu věci dělat co si jen zachce? Chceme vůbec ve hře mít výrazné telegrafování vzhledem k tomu, že neodpovídá realitě? Rozlousknutí tohoto problému si vyžaduje kompromisy.
\bigbreak
V neposlední řadě, plná kontrola nad mečem nemusí vhodně ladit s \textbf{narativní} složkou hry. V mnoha akčních hrách prezentuje příběh hráčovu postavu jako zkušeného šermíře, ne-li přímo virtuosa meče, jakému se v herním světě vyrovná jen hrstka dalších. Problémem však je, že většinový hráč srovnatelně zkušeným šermířem s největší pravděpodobností není. Je-li šermířská dovednost postavy fakticky přímo úměrná hráčově dovednosti, může tím tedy docházet k narativní disonanci. Zahrnutím takového bojového systému se designérům, kteří mají narativ na starosti, citelně komplikuje práce. 
\bigbreak
Finálně, pokud jsme všechny tyto problémy vyřešili, čeká nás implementace systému. Hráči jsme nabídli v podstatě neomezené možnosti, nabízí se tedy otázka, jakým způsobem se hodláme ujistit, že možnosti skutečně neomezené jsou a nevyskytly se nějaké \textbf{neočekávané okrajové případy}, které nebyly ošetřeny a povedou k nedefinovanému chování? Jednou z cest je přeci jen kontrolu nad mečem omezit a smrsknout ji na stále velkou, avšak \textbf{konečnou množinu stavů}, které jsou tvůrci schopni projít jednu po druhé a ujistit se, že každý funguje, jak zamýšleli. Jedinou další alternativou se zdá být \textbf{plně realistická fyzikální simulace meče} - problémům, které tato cesta přináší, je věnována převážná část této práce.  
\bigbreak
Vidíme, že poskytnout hráči přímou kontrolu nad mečem, je nevděčný úkol vyžadující mnoho kompromisů. V historii najdeme pouze několik her, které se s ním pokusily popasovat, každá svým unikátním způsobem. Na ty se nyní podíváme blíže.

\subsection{Die by the Sword}

Průkopníkem v této oblasti je Die by the Sword \cite{DieByTheSword}, hra vydaná kalifornskou společností Treyarch Corporation roku 1998. 

Jde o jedinou hru, kterou si zde představíme, jež se nevydala cestou omezení množiny stavů meče na konečné množství, ale nabídla hráči skutečně \textbf{jemnou, spojitou kontrolu}. Ta byla umožněna zapojením jednoduchého \textbf{realistického fyzikálního subsystému}, což byl na tehdejší dobu úkaz vskutku nevídaný vzhledem k hardwarovým omezením, v jejichž rámci bylo nutno operovat.

\textbf{Backend hry} umožňoval \textbf{plně volný pohyb} všech objektů třímaných herní postavou (postava mohla např. v jedné ruce držet meč a ve druhé štít či druhý meč), přičemž tělo postavy bylo \textbf{procedurálně animováno} (další na tuto dobu vzácný úkaz), aby poloha jejích rukou s polohou objektů korespondovala.

Tento přístup k implementaci přinesl některé \textbf{zajímavé důsledky}. Hráč se nemusel učit používat dedikovanou mechaniku pro blokování nepřítelských úderů, stačilo meč napolohovat tak, aby tomu nepřátelskému stál v cestě. Podobným způsobem zde byl nepřeberný prostor pro další zajímavé techniky, vyvíjené a zdokonalované plně hráčskou komunitou. Dále, díky obecnosti systému nebyl problém, mít nepřátele všech představitelných tvarů a velikostí - dokud měl nepřítel hmotné tělo, mohli jste ho mečem trefit; dokud měl nějakou končetinu, mohl sám držet zbraň. Spojením se systémem procedurální animace, nestvůrám bylo možné velmi efektně odsekávat kusy těla a kochat se jejich zvýšenou nemotorností\footnote{případně rovnou neschopností hráče zranit, pokud byla odseknuta končetina třímající zbraň}. K dovršení dojmu vás hra nechala chopit se odseknuté končetiny jako zbraně, stále při stejném ovládání, jaké vám nabídl meč.

Kamenem úrazu pro hru však nesporně bylo \textbf{ovládání}. Tvůrci hry si byli zřejmě vědomi, o jak náročnou výzvu vyžadující dlouhodobý výzkum v případě této hry jde, a nabídli hned několik velmi odlišných metod vstupu pro udání pohybu zbraně:
\begin{itemize}
    \item \textbf{Numerická klávesnice}\footnote{V době vzniku hry nebyl grafický operační systém ovládaný myší běžnou součástí výbavy každého hráče. Od her se očekávalo, že budou podporovat možnost ovládání čistě pomocí klávesnice.} - Obdélník číselných kláves symbolizoval obdélník obrazovky, stisknutí dvou kláves (poloh na obrazovce) v rychlém sledu znamenalo švih mečem z jedné polohy na druhou.
    \item \textbf{Myš} - Myší bylo možné udávat pohyb spojitým způsobem. Poloha kurzoru na obrazovce se mapovala na polohu v herním světě, kam se má švihem meč přesunout. 
    \item \textbf{Nahrávky} - Ve hře byl obsažen dedikovaný nástroj umožňující na časové ose, pomocí klíčových snímků - způsobem klasicky používaným v nástrojích pro tvorbu animací - vytvářet a editovat vlastní zcela libovolné pohyby zbraně. Hotové nahrávce bylo následně možné přiřadit klávesu, jejímž stisknutím ji postave ve hře přehrála. 
\end{itemize}   
Přes značné úsilí, které tvůrci této oblasti věnovali, stalo se ovládání hlavním předmětem hráčské kritiky. Mnoho hráčů si stěžovalo, že je hra jeho vinou nadmíru těžká, případně že si zkrátka \textbf{nehodlají na neotřelý způsob ovládání zvykat}. Přestože hra byla značně oceněna kritiky, mezi širokou herní veřejností se jí výrazných ohlasů nedostalo. Kromě datadisku Limb from Limb se již na její odkaz žádná další hra nepokusila navázat, pokud je autorovi této práce známo.

\subsection{Mount\&Blade}

Kultovní klasikou je Mount\&Blade. První verzi této hry vydalo turecké studio TaleWorlds v roce 2008, později se hra pro velkou oblíbenost mezi hráči dočkala obsáhlého datadisku Warband a velmi nedávno i druhého dílu Bannerlord, který mnoho mechanik hrubě nastíněných v původní hře dále rozvádí a prohlubuje. My se zde budeme věnovat prvnímu dílu, protože ten je v kombinaci s datadiskem Warband i v dnešní době mezi hráči velmi oblíbený a materiálů pro analýzu nám poskytne zcela dostatečné množství.




\clearpage
(stručné představení hry)

(popis bojového systému)

(jaké výhody z tohoto plynou)

(jaké výhody se nepodařilo získat)

\subsection{Kingdom Come: Deliverance}

Kingdom Come: Deliverance bylo vydáno roku 2018 českými Warhorse Studios. 

\clearpage
(stručné představení hry)

(popis bojového systému)

(jaké výhody z tohoto plynou)

(jaké výhody se nepodařilo získat)


\section{Alternativní ovladače}

(představení - ovladače co mohou být pro náš účel ovládání meče víc vhodné - existují ovladače co se postupně dostávají do mainstreamu, co jsou state-of-the art ve svém obskurním niche oboru, i naprosto experimentální ovladače co jsou teprve předmětem výzkumu)

\subsection{Virtuální realita}

(jaké má výhody)

(příklad hry kde se s ním ovládá meč) (swords\&sorcery, gorn, beatsaber)

(nevýhody)
- disonance "meč narazil do něčeho a není tam kde ovladač"

(proč se jí nebudeme při výzkumu zabývat)
- pořád není mainstream
- (zjistit) uzavřenost platforem?


\subsection{Pohybové senzory ve smartphonu}

(výhody)
- dostupnost
- otevřenost platformy, programátorská flexibilita (hackovatelnost)

(nevýhody)
- fragmentace mobilního trhu, nekonzistence výstupu senzorů mezi různými mobily
- způsob připojení (ne každej desktop má bluetooth nebo wifi, kabel je nepraktickej; celkově mobily asi maj problém s nízkolatenčním přenosem po místní síti)
- najít nějakej článek

\subsection{Obskurní periferie}

3d myš
- extrémně drahá 
- stavěná pro 3d modelováni
- nemá 6 ale jenom 3 stupně volnosti


\section{Shrnutí}

%%% Fiktivní kapitola s ukázkami citací

\chapter{Odkazy na literaturu}

Odkazy na literaturu vytváříme nejlépe pomocí příkazů
\verb|\citet|, \verb|\citep| atp.
(viz {\LaTeX}ový balíček \textsf{natbib}) a~následného použití
Bib{\TeX}u. V~matematickém textu obvykle odkazujeme stylem \uv{Jméno
autora/autorů (rok vydání)}, resp. \uv{Jméno autora/autorů [číslo
odkazu]}. V~českém/slovenském textu je potřeba se navíc vypořádat
s~nutností skloňovat jméno autora, respektive přechylovat jméno
autorky. Je potřeba mít na paměti, že standardní příkazy
\verb|\citet|, \verb|\citep|
produkují referenci se jménem autora/autorů v~prvním pádě a~jména
autorek jsou nepřechýlena.

Pokud nepoužíváme bib\TeX{}, řídíme se normou ISO 690 a zvyklostmi
oboru.

Jména časopisů lze uvádět zkráceně, ale pouze v~kodifikované podobě.

\section{Několik ukázek}

Mezi nejvíce citované statistické články patří práce Kaplana a~Meiera a~Coxe
\citep{KaplanMeier58, Cox72}. \citet{Student08} napsal článek o~t-testu.

Prof. Anděl je autorem učebnice matematické statistiky
\citep[viz][]{Andel98}. Teorii odhadu se věnuje práce
\citet{LehmannCasella98}. V~případě odkazů na specifickou informaci
(definice, důkaz, \dots) uvedenou v~knize bývá užitečné uvést
specificky číslo kapitoly, číslo věty atp. obsahující požadovanou
informaci, např. viz \citet[Věta 4.22]{Andel07} nebo \citep[viz][Věta
4.22]{Andel07}.

Mnoho článků je výsledkem spolupráce celé řady osob. Při odkazování
v~textu na článek se třemi autory obvykle při prvním výskytu uvedeme
plný seznam: \citet*{DempsterLairdRubin77} představili koncept EM
algoritmu. Respektive: Koncept EM algoritmu byl představen v~práci
Dempstera, Lairdové a~Rubina \citep*{DempsterLairdRubin77}. Při každém
dalším výskytu již používáme zkrácenou verzi:
\citet{DempsterLairdRubin77} nabízejí též několik příkladů použití EM
algoritmu. Respektive: Několik příkladů použití EM algoritmu lze
nalézt též v~práci Dempstera a~kol. \citep{DempsterLairdRubin77}.

U~článku s~více než třemi autory odkazujeme vždy zkrácenou formou:
První výsledky projektu ACCEPT jsou uvedeny v~práci Genbergové a~kol.
\citep{Genberget08}. V~textu \emph{nenapíšeme}: První výsledky
projektu ACCEPT jsou uvedeny v~práci \citet*{Genberget08}.

\chapter{Technické základy}

V této kapitole čtenáře seznámíme s herním enginem Unity \cite{Unity} do míry potřebné k pochopení zbytku práce. Stručně nastíníme architekturu, kterou pro tvorbu herních systémů tento engine podporuje, a zmíníme některé jeho vestavěné knihovny - především se podrobně zaměříme na integrovaný fyzikální subsystém.


\section{Herní engine Unity} \label{unityEngineIntroSection}

První verze enginu Unity byla vydána v roce 2005. Od té doby bylo vydáno mnoho verzí, přidáno mnoho funkcionalit a mnoho dalších prohlášeno za zastaralé. Tvůrcům se dařilo držet krok s dobou, v dnešní době jde o velmi vyspělý a moderní engine, který uživatelům nabízí \textbf{široký ekosystém nástrojů}, ať už jde o integrované komponenty, nebo balíčky třetích stran. 

Jednou z jeho velkých sil je \textbf{všestrannost} - zvládá stejně dobře 2D i 3D hry, výsledný spustitelný produkt lze vyexportovat pro širokou škálu platforem - desktopové, mobilní, konzole i web.

Pro programování herní logiky je zde podporován jazyk C\#. Jako runtime je dosud používáno postarší \textbf{Mono}, které umožňuje nanejvýše \textbf{C\# verze 9} s pár drobnostmi nepodporovanými, avšak stále jde o vcelku moderní jazyk, který programátorovi umožňuje pohodlnou práci. Interakce s knihovnami třetích stran, které závisejí na některých součástech moderního frameworku .NET, může činit problémy, avšak typicky nejde o nepřekonatelnou překážku.

Jde o \textbf{komerční produkt s uzavřenou licencí}, avšak cenová politika je více než přívětivá. Pro nekomerční projekty a projekty nepřesahující určitou hranici celkového ročního zisku je \textbf{použití enginu zcela zdarma}. Právní úprava nijak nebrání vydávání samostatných projektů v enginu tvořených pod opensource licencí. Dokumentace není celkově na špatné úrovni a těch pár míst, kde pokulhává, vynahrazují velmi živá komunitní fóra.
\bigbreak
V této kapitole si představíme základní principy použití enginu v rozsahu nutném pro vysvětlení zbytku práce. Jejich kompletnější popis čtenář najde v oficielní dokumentaci \cite{Unity}. 

\subsection{Herní objekty a komponenty}
\begin{itemize}
    \item scéna - hierarchie gameobjectů
    \item gameobject - kontejner do kterého se naházejí komponenty
    \item komponenta - modul poskytující objektu nějaký kousek funkcionality; builtin, nebo to je uživatelsky definovaná třída; může obsahovat nějaká data (nastavitelná v editoru) a nějakou herní logiku (která příp. zahrnuje komunikaci s ostatními komponentami a gameobjecty)
    \item serializace - z editoru se dají editovat public/[SerializeField] fieldy ; nekompatibilní s properties, interfacama a jinýma datovýma strukturama než pole/list - omezující když chce člověk používat klasické dobré OOP practices  
    \item z herního objektu se dá vytvořit prefab
\end{itemize}


\subsection{Posílání zpráv}
\begin{itemize}
    \item srovnání VS UnityEvent VS C\# delegáti 
\end{itemize}

\subsection{Herní smyčka}
\begin{itemize}
    \item zpráva Update, deltaTime
    \item všechno probíhá sekvenčně na main threadu
\end{itemize}

\subsection{Základní knihovny}
\begin{itemize}
    \item (vše pouze zmínit že existuje)
    \item matematická primitiva - Vector3, Ray, matice, apod.
    \item instanciace gameobjectů a komponent za runtimu
    \item Vykreslování grafiky
    \item NavMesh
    \item extrémně moc dalších věcí
\end{itemize}


\section{Fyzikální subsystém enginu Unity}

\subsection{Představení}
\begin{itemize}
    \item zjednodušená dynamika a kinematika pevného tělesa
    \item je to vlastní Uniťácké API postavené nad backendem nVidia PhysX, ale poskytuje jenom podmnožinu kompletní PhysX funkcionality
    \item uvést reference na Unity a na PhysX dokumentaci\footnote{Chci mít odděleně referenci na Unity a na Unity-docs, nebo jenom Unity a nechat ho aby url neukazovalo na main page ale na docs?}
    \item zmínit že Unity má taky 2D Physics, co jsou úplně odděleně a za backend maj Box2D 
    \item vede si vlastní reprezentaci celé scény, při počítání vychází jenom z ní, před a po každém kroku se provede synchronizace s herním světem (updatnou transformy apod.)
    \item FixedUpdate
\end{itemize}


\subsection{Rigidbody}
\begin{itemize}
    \item non-kinematic VS kinematic
    \item properties (mass, center-of-mass+inertia-tensor, (angular) drag, axis-lock,...)
    \item jak se chová vůči transform-hierarchii
          \begin{itemize}
            \item fyzika používá globální souřadnice objektů
            \item collidery v potomcích patří otcovskému rigidbody
            \item rigidbody uvnitř rigidbody -> undefined behavior
          \end{itemize}
\end{itemize}

\subsection{Colliders}
\begin{itemize}
    \item zmínit že to je nejvíc základní a používaná funkcionalita physics systému, kvůli který ho spousta her includuje i když nepoužívají žádný non-kinematic rigidbodies a tak.  
    \item primitivní tvary, jejich skládání, nefunkčnost obecné maticové deformace (vždy zachovává kvádrovost, kulovost, kapsulovost)
    \item NEodpovídá(!) přesně tvaru meshe a v tom je pointa, ale když se umístěj šikovně, nijak to nekazí hráčský zážitek
    \item (letmo zmínit) mesh collidery existují - ale je u nich dost omezení, mají smysl jenom na terén apod.
    \item statické vs dynamické collidery (úplně samostatné / kinematic / non-kinematic) - co s čím koliduje a jak se chovaj když se pohybujou (pohyblivé collidery by měly být na nějakém (klidně kinematickém) rigidbody)
    \item trigger collidery
    \item Layers - co s čím koliduje; slouží i pro jinou funkcionalitu naprosto nesouvisející s fyzikou(!) ( další problémy rozebereme v \ref{swordCollisionsSection})
    \item zpráva OnCollisionXx / OnTriggerXx
    \item discrete VS continuous collision detection (jenom že existuje, podrobněji v \ref{swordCollisionsSection})
\end{itemize}

\subsection{Joints}
\begin{itemize}
    \item (velmi stručně) co je constraint solver a jaký používá Unity
    \item strohý výčet jmen jointů a že všechny jsou to speciální případ Configurable Jointu (neboli D6 jointu jak ho zná PhysX)
    \item představení Configurable Jointu
    \item anchor, connectedAnchor
    \item locked/limited/free pro každý degree of freedom
    \item primary, secondary, tertiary axis, vůči čemu jsou rotace relativní (podrobněji rozebereme v \ref{howToSetJointsTargetRotationSection})
    \item positional\&angular drive - síly (spring VS damper), cílová pozice (nestabilní, alternativa nastavení connectedAnchoru) a rotace, RotationDriveMode (Slerp VS X+YZ - jenom letmo zmínit, "Slerp se choval pěkně a tak jsme ho vybrali"), Target (angular) velocity
    \item massScale VS connectedMassScale
    \item stabilita, flag EnablePreprocessing, projekce (projection mode+distance+angle)
\end{itemize}


\subsection{Stabilita simulace}
\begin{itemize}
    \item všechno to je jenom hodně hrubá aproximace, priorita je aby to běželo rychle a dalo se rozběhnout na GPU
    \item postup simulace po krocích, tunelování
    \item přesnost IEEE floatu - hodně malé/velké objekty, zaokrouhlovací chyba
    \item constraint solver, nesplnitelné constrainty, iterativní metody
    \item to, že collidery neodpovídají meshům, taky nepřidává přesnosti
    \item ladění - parametry fixedFrameDuration, iterationCount, který constraint solver se používá apod.
    \item zmínit vestavěný Physics Debugger
\end{itemize}


\section{Shrnutí}

Unity je velmi vyspělý engine, který za sebou má desetiletí vývoje a s jeho pomocí bylo vytvořeno mnoho komerčně úspěšných her. Pro svou velmi přívětivou cenovou politiku a možnost kompilovat hry pro širokou škálu platforem od konzolí po mobilní zařízení se stal \textbf{oblíbeným nástrojem začínajících vývojářů a malých týmů}. Uživateli poskytuje mnoho nástrojů a knihoven, od editoru scény po subsystém pro navigaci entit v herním světě. Libovolnou další funkcionalitu lze přidat pomocí skriptů v jazyce C\#. 

Herní scéna je hierarchií objektů - každý z nich je definován svým předkem, jménem, transformační maticí a listem komponent. Komponenty jsou základním stavebním prvkem, kterým dodáváme objektům jejich vlastnosti - např. schopnost mít vizuální reprezentaci, schopnost kolidovat s okolními objekty. Unity poskytuje mnoho předpřipravených komponent, uživatel může definovat vlastní jako skripty dědící z třídy MonoBehavior.

Jednou ze zajímavých skupin komponent jsou ty poskytnuté vestavěným \textbf{fyzikálním sybsystémem}. Jejich prostřednictvím je uživateli poskytnuto rozhraní odpovídající velké podmnožině systému Nvidia PhysX \cite{PhysX}. 

Základem je komponenta Rigidbody - jejím přidáním se z objektu stává fyzikální entita, na kterou lze působit silami. Chceme-li fyzikálnímu objektu umožnit kolidovat s dalšími objekty, přidáme mu jednu nebo více komponent typu Collider. Pro vytváření složitějších vztahů a závislostí mezi jednotlivými fyzikálními tělesy slouží Jointy - s jejich pomocí můžeme vytvořit například dveře otevíratelné v pantech či realisticky se chovající kovový řetěz.

Fyzikální simulace je velmi mocný nástroj, nevytváří však věrný obraz reality. Na mnoha místech narážíme na numerickou stabilitu a také zkrátka na fakt, že hlavní prioritou nikdy nebyla dokonalá přesnost, nýbrž schopnost vyhodnocení v reálném čase.

%%% Fiktivní kapitola s instrukcemi k PDF/A

\chapter{Formát PDF/A}

Opatření rektora č. 13/2017 určuje, že elektronická podoba závěrečných
prací musí být odevzdávána ve formátu PDF/A úrovně 1a nebo 2u. To jsou
profily formátu PDF určující, jaké vlastnosti PDF je povoleno používat,
aby byly dokumenty vhodné k~dlouhodobé archivaci a dalšímu automatickému
zpracování. Dále se budeme zabývat úrovní 2u, kterou sázíme \TeX{}em.

Mezi nejdůležitější požadavky PDF/A-2u patří:

\begin{itemize}

\item Všechny fonty musí být zabudovány uvnitř dokumentu. Nejsou přípustné
odkazy na externí fonty (ani na \uv{systémové}, jako je Helvetica nebo Times).

\item Fonty musí obsahovat tabulku ToUnicode, která definuje převod z~kódování
znaků použitého uvnitř fontu to Unicode. Díky tomu je možné z~dokumentu
spolehlivě extrahovat text.

\item Dokument musí obsahovat metadata ve formátu XMP a je-li barevný,
pak také formální specifikaci barevného prostoru.

\end{itemize}

Tato šablona používá balíček {\tt pdfx,} který umí \LaTeX{} nastavit tak,
aby požadavky PDF/A splňoval. Metadata v~XMP se generují automaticky podle
informací v~souboru {\tt prace.xmpdata} (na vygenerovaný soubor se můžete
podívat v~{\tt pdfa.xmpi}).

Validitu PDF/A můžete zkontrolovat pomocí nástroje VeraPDF, který je
k~dispozici na \url{http://verapdf.org/}.

Pokud soubor nebude validní, mezi obvyklé příčiny patří používání méně
obvyklých fontů (které se vkládají pouze v~bitmapové podobě a/nebo bez
unicodových tabulek) a vkládání obrázků v~PDF, které samy o~sobě standard
PDF/A nesplňují.

Další postřehy o~práci s~PDF/A najdete na \url{http://mj.ucw.cz/vyuka/bc/pdfaq.html}.

\chapter{Zhodnocení}
gfdew

\section{Shrnutí}
\chapter{Vyhodnocení}


\section{Shrnutí}

\chapter*{Závěr}
\addcontentsline{toc}{chapter}{Závěr}


%%% Seznam použité literatury
%%% Seznam použité literatury (bibliografie)
%%%
%%% Pro vytváření bibliografie používáme bibTeX. Ten zpracovává
%%% citace v textu (např. makro \cite{...}) a vyhledává k nim literaturu
%%% v souboru literatura.bib.
%%%
%%% Příkaz \bibliographystyle určuje, jakým stylem budou citovány odkazy
%%% v textu. V závorce je název zvoleného souboru .bst. Styly plainnat
%%% a unsrt jsou standardní součástí latexových distribucí. Styl czplainnat
%%% je dodáván s touto šablonou a bibTeX ho hledá v aktuálním adresáři.

%\bibliographystyle{czplainnat}    %% Autor (rok) s českými spojkami
% \bibliographystyle{plainnat}    %% Autor (rok) s anglickými spojkami
% \bibliographystyle{unsrt}       %% [číslo]

\renewcommand{\bibname}{Seznam použité literatury}

%%% Vytvoření seznamu literatury. Pozor, pokud jste necitovali ani jednu
%%% položku, seznam se automaticky vynechá.

\printbibliography
\addcontentsline{toc}{chapter}{Seznam použité literatury}
%\bibliography{literatura}

%%% Kdybyste chtěli bibliografii vytvářet ručně (bez bibTeXu), lze to udělat
%%% následovně. V takovém případě se řiďte normou ISO 690 a zvyklostmi v oboru.

% \begin{thebibliography}{99}
%
% \bibitem{lamport94}
%   {\sc Lamport,} Leslie.
%   \emph{\LaTeX: A Document Preparation System}.
%   2. vydání.
%   Massachusetts: Addison Wesley, 1994.
%   ISBN 0-201-52983-1.
%
% \end{thebibliography}


%%% Obrázky v bakalářské práci
%%% (pokud jich je malé množství, obvykle není třeba seznam uvádět)
\listoffigures

%%% Tabulky v bakalářské práci (opět nemusí být nutné uvádět)
%%% U matematických prací může být lepší přemístit seznam tabulek na začátek práce.
\listoftables

%%% Použité zkratky v bakalářské práci (opět nemusí být nutné uvádět)
%%% U matematických prací může být lepší přemístit seznam zkratek na začátek práce.
\printacronyms[heading=chapter*, name={Seznam zkratek}]
\addcontentsline{toc}{chapter}{Seznam zkratek}

%%% Přílohy k bakalářské práci, existují-li. Každá příloha musí být alespoň jednou
%%% odkazována z vlastního textu práce. Přílohy se číslují.
%%%
%%% Do tištěné verze se spíše hodí přílohy, které lze číst a prohlížet (dodatečné
%%% tabulky a grafy, různé textové doplňky, ukázky výstupů z počítačových programů,
%%% apod.). Do elektronické verze se hodí přílohy, které budou spíše používány
%%% v elektronické podobě než čteny (zdrojové kódy programů, datové soubory,
%%% interaktivní grafy apod.). Elektronické přílohy se nahrávají do SISu a lze
%%% je také do práce vložit na CD/DVD. Povolené formáty souborů specifikuje
%%% opatření rektora č. 72/2017.
\appendix
\chapter{Přílohy}

\section{První příloha}

\openright
\end{document}
